% Figure: current-voltage curve of a resonant-tunnelling diode. As the emitter
% level aligns with the quantum-well resonance the current peaks; pushing the
% bias past alignment cuts off the resonant channel and the current falls, the
% negative-differential-resistance region that makes the device an oscillator.
% All points entered explicitly; no runtime computation.
\begin{figure}[htbp]
\centering
\begin{tikzpicture}[font=\scriptsize]
\begin{axis}[
    width=10cm, height=6.2cm,
    xlabel={bias voltage $V$ (\si{\volt})},
    ylabel={current $I$ (mA)},
    xmin=0, xmax=1.0,
    ymin=0, ymax=6,
    xtick={0,0.2,0.4,0.6,0.8,1.0},
    ytick={0,1,2,3,4,5,6},
    grid=both,
    grid style={enggray!25},
    tick label style={font=\tiny},
    label style={font=\scriptsize},
    axis line style={enggray},
    every axis plot/.append style={very thick},
]
% resonant rise to peak, NDR valley, then thermionic rise
\addplot[engblue, mark=*, mark size=1.0pt] coordinates {
  (0.00,0.00) (0.05,0.15) (0.10,0.45) (0.15,1.10) (0.20,2.20)
  (0.25,3.60) (0.30,4.80) (0.34,5.20) (0.36,5.05) (0.40,4.10)
  (0.45,2.60) (0.50,1.40) (0.55,0.90) (0.60,0.80) (0.65,0.85)
  (0.70,1.05) (0.80,1.90) (0.90,3.40) (1.00,5.20)
};
% mark peak
\addplot[engred, only marks, mark=o, mark size=3pt] coordinates {(0.34,5.20)};
\node[engred, anchor=south, font=\tiny] at (axis cs:0.34,5.30) {peak};
% mark valley
\addplot[engorange, only marks, mark=square*, mark size=2.2pt]
  coordinates {(0.60,0.80)};
\node[engorange, anchor=north west, font=\tiny] at (axis cs:0.61,0.85)
  {valley};
% annotate NDR region
\node[enggray, font=\tiny, align=center] at (axis cs:0.47,3.4)
  {negative\\differential\\resistance};
\draw[enggray, ->] (axis cs:0.47,2.9) -- (axis cs:0.44,2.6);
\end{axis}
\end{tikzpicture}
\caption[Resonant-tunnelling-diode current-voltage curve]{Current against bias
for a double-barrier resonant-tunnelling diode. At low bias the emitter Fermi
sea aligns with the bound state of the quantum well between the two barriers and
electrons tunnel resonantly, the current rising to a peak. Raising the bias past
alignment lifts the emitter above the well level, closes the resonant channel,
and the current falls even as the voltage rises, the region of negative
differential resistance between the peak and the valley. Beyond the valley,
ordinary thermionic and non-resonant tunnelling take over and the current climbs
again. The falling region delivers gain and is the basis of the device's use as
a high-frequency oscillator; the peak-to-valley current ratio is its figure of
merit.}
\label{fig:vol04ch12:rtd-iv}
\end{figure}
